\subsection{从语法到算子：转移算子、C*-代数与 \texorpdfstring{$\zeta$}{zeta}（精确口径）}\label{subsec:operator-zeta-interface}

本小节把“语法—谱—$\zeta$”的联系写成可评估的算子论陈述，并明确本文真正使用的对象是什么。核心点是：在黄金均值 SFT 上，Perron--Frobenius（有限矩阵）与 Ruelle 转移算子（无穷维算子）在柱函数空间上一致；相应的 C*-代数框架（如 Cuntz--Krieger 代数/群胚 C*-代数）提供了把柱投影、最大熵态（Parry 测度）与谱/$\zeta$ 统一表述的标准工具。

\subsubsection{转移算子与 \texorpdfstring{$\zeta$}{zeta} 的 Fredholm 行列式口径}

令 $X_A^+\subset\{0,1\}^{\NN}$ 为由邻接矩阵 $A$ 定义的一侧 SFT（黄金均值约束），令 $\sigma:X_A^+\to X_A^+$ 为左移位。定义零势（计数势）下的 Ruelle 转移算子（对有界函数）
\[
(\mathcal{L}f)(x):=\sum_{\sigma y=x} f(y).
\]
在 SFT 情形，该求和是有限的。令 $\mathcal{C}_m$ 为由长度-$m$ 柱集指示函数张成的有限维子空间，则 $\mathcal{L}$ 在 $\mathcal{C}_1$ 上等价于矩阵 $A^{\mathsf T}$ 的作用，并由柱函数闭包可追踪到所有有限 $m$。

\begin{definition}[Fredholm 行列式（可审计口径）]\label{def:fredholm-determinant}
若 $T$ 为迹类算子（在本文使用的有限维/有限秩情形自动成立），定义其 Fredholm 行列式为
\[
\det(I-zT):=\exp\left(-\sum_{n\ge 1}\frac{\mathrm{Tr}(T^n)}{n}z^n\right),
\quad |z|\ \text{充分小}.
\]
在有限维情形该定义与通常矩阵行列式一致，并由解析延拓得到有理函数恒等式（参见 \cite{Ruelle1976ZetaExpanding,ParryPollicott1990Zeta}）。
\end{definition}

对有限型子移位（SFT），动力学 $\zeta$ 的标准恒等式为
\[
\zeta_\sigma(z)=\frac{1}{\det(I-zA)},
\]
并可由周期点计数与 Fredholm 行列式口径一致地给出（参见 \cite{LindMarcus1995,Ruelle1976ZetaExpanding,Ruelle1978Thermodynamic,ParryPollicott1990Zeta}）。

\subsubsection{带“复杂度势”的二变量 \texorpdfstring{$\zeta$}{zeta} 指纹：Euler 乘积与迹序列的严格对齐}\label{subsec:zeta-two-variable-kappa}

本节给出一个对本文“prime-like（Euler 乘积/显式公式）侧”与“谱/迹（trace）侧”都兼容、并可直接用于数值审计的统一对象：\textbf{带势的两变量 \(\zeta_K(z,u)\)}。其中 \(z\) 计时间步数（迹的幂次），\(u\) 计一类局部可加的“复杂度/串行代价势” \(\kappa\)（例如不可并行的 delay/carry 事件数），从而把“时间代价”以乘法权重的形式嵌入 Euler 乘积的 primitive 因子。

\begin{definition}[带势的两变量 \texorpdfstring{$\zeta_K(z,u)$}{zetaK(z,u)}]\label{def:zetaK-zu}
给定一个有限有向图（或有限状态同步系统）\(\Gamma_K\)，并对每条边 \(e\) 赋一个非负整数势
\[
\kappa(e)\in\ZZ_{\ge 0}.
\]
定义带权邻接矩阵族 \(M_K(u)\in \Mat_{|V|\times |V|}(\ZZ[u])\) 为
\[
\bigl(M_K(u)\bigr)_{ij}:=\sum_{e:i\to j} u^{\kappa(e)}.
\]
定义两变量 \(\zeta\) 为形式幂级数（或在 \(|z|\) 充分小时的解析函数）
\begin{equation}\label{eq:zetaK-trace-exp}
\zeta_K(z,u)
:=\exp\!\Bigl(\sum_{n\ge 1}\frac{a_n(u)}{n}z^n\Bigr),
\qquad
a_n(u):=\Tr\bigl(M_K(u)^n\bigr),
\end{equation}
从而在有限维情形严格有
\begin{equation}\label{eq:zetaK-det}
\zeta_K(z,u)=\frac{1}{\det(I-zM_K(u))}.
\end{equation}
\end{definition}

\begin{remark}[Euler 乘积与 primitive 分层谱]\label{rem:zetaK-euler-product}
记 \(\mathcal{P}_K\) 为 \(\Gamma_K\) 的 primitive 有向闭轨道集合；对 \(\gamma\in\mathcal{P}_K\) 记其长度为 \(|\gamma|\)，并定义轨道势为边势的轨道和
\[
\kappa(\gamma):=\sum_{e\in\gamma}\kappa(e).
\]
则由标准的“周期点/闭路”展开（参见 \cite{ParryPollicott1990Zeta,LindMarcus1995}）得到严格的形式乘积恒等式
\begin{equation}\label{eq:zetaK-euler}
\zeta_K(z,u)=\prod_{\gamma\in\mathcal{P}_K}\bigl(1-z^{|\gamma|}u^{\kappa(\gamma)}\bigr)^{-1}.
\end{equation}
引入 primitive 分层谱多项式
\begin{equation}\label{eq:zetaK-Bn}
B_n(u):=\sum_{\gamma\in\mathcal{P}_K,\ |\gamma|=n} u^{\kappa(\gamma)}\in\ZZ[u],
\end{equation}
则 \(\{B_n(u)\}_{n\ge 1}\) 与迹序列 \(\{a_n(u)\}_{n\ge 1}\) 之间存在严格的 Möbius/Witt 互换（见下命题）。
\end{remark}

\begin{proposition}[Witt/Möbius 反演：由迹序列提取 primitive 分层谱]\label{prop:zetaK-mobius-primitive}
沿用定义 \ref{def:zetaK-zu}--\ref{rem:zetaK-euler-product} 的记号。则对任意 \(n\ge 1\)：
\begin{equation}\label{eq:zetaK-trace-from-primitive}
a_n(u)=\sum_{d\mid n} d\,B_d\!\bigl(u^{n/d}\bigr),
\end{equation}
并且由 Möbius 反演可严格恢复
\begin{equation}\label{eq:zetaK-primitive-from-trace}
\boxed{\;
B_n(u)=\frac{1}{n}\sum_{d\mid n}\mu(d)\,a_{n/d}\!\bigl(u^{d}\bigr)\; }.
\end{equation}
\end{proposition}

\begin{proof}
把 Euler 乘积 \eqref{eq:zetaK-euler} 取对数并展开为几何级数，得
\[
\log\zeta_K(z,u)
=\sum_{\gamma\in\mathcal{P}_K}\sum_{k\ge 1}\frac{1}{k}z^{k|\gamma|}u^{k\kappa(\gamma)}.
\]
将右侧按 \(n=k|\gamma|\) 分组，并注意每个 primitive 轨道 \(\gamma\) 的 \(k\) 次幂产生 \( |\gamma| \) 条不同起点的长度 \(n\) 闭路，从而长度 \(n\) 的迹计数满足 \eqref{eq:zetaK-trace-from-primitive}。再对 \eqref{eq:zetaK-trace-from-primitive} 作 Möbius 反演即得 \eqref{eq:zetaK-primitive-from-trace}。证毕。
\end{proof}

\begin{remark}[严格交换的“迹 $\to$ primitive $\to$ Euler”桥包（交换图）]\label{rem:zetaK-commuting-diagram}
命题 \ref{prop:zetaK-mobius-primitive} 与定义 \ref{def:zetaK-zu}（行列式口径 \eqref{eq:zetaK-det}）合在一起，给出本文反复调用的\textbf{严格交换}桥包：对任意有限核/有限矩阵 \(T\)（例如固定 \(u\) 后的 \(T=M_K(u)\)），令
\[
a_n(T):=\Tr(T^n)\quad(n\ge 1),
\qquad
p_n(T):=\frac{1}{n}\sum_{d\mid n}\mu(d)\,a_{n/d}(T),
\]
并定义 Euler 变换
\[
\mathrm{Euler}\bigl((p_n)_{n\ge 1}\bigr)(z):=\prod_{n\ge 1}(1-z^n)^{-p_n}.
\]
则对每个 \(T\) 都有逐项恒等式
\[
\boxed{\;
\zeta_T(z):=\frac{1}{\det(I-zT)}
=\exp\!\Bigl(\sum_{n\ge 1}\frac{a_n(T)}{n}z^n\Bigr)
=\prod_{n\ge 1}(1-z^n)^{-p_n(T)}\; }.
\]
换言之，把“迹序列（ghost）$\to$ Möbius/Witt 反演（primitive）$\to$ Euler 乘积（prime-like）”压成最小结构，可写为如下严格交换图：
\[
\begin{CD}
T @>\mathrm{TrSeq}>> (a_n(T))_{n\ge 1} @>\mathrm{Mob}>> (p_n(T))_{n\ge 1} @>\mathrm{Euler}>> \prod_{n\ge 1}(1-z^n)^{-p_n(T)}\\
@V\mathrm{ZetaDet}VV @. @. @VV=V\\
\det(I-zT)^{-1} @. @. \det(I-zT)^{-1}
\end{CD}
\]
其中 \(\mathrm{TrSeq}(T)=(a_n(T))\)、\(\mathrm{Mob}(a)=p\) 为 Möbius 反演，而 \(\mathrm{ZetaDet}(T)=\det(I-zT)^{-1}\)。
\end{remark}

\begin{remark}[Fourier--Witt 双幽灵桥（FW-Bridge）]\label{rem:fw-bridge}
注 \ref{rem:zetaK-commuting-diagram} 给出时间/谱侧的“迹 ghost \(\to\) primitive \(\to\) Euler”严格交换桥包（Witt 侧）。
另一方面，Fold 支路在有限循环环 \(G_m=\ZZ/F_{m+2}\ZZ\) 上还有一条完全独立、同样严格交换的“卷积 \(\leftrightarrow\) 点乘 + Parseval 能量”桥包（Fourier 侧；见注 \ref{rem:fold-fourier-parseval-bridge}）。
两条桥共享同一个可计算的“幽灵标尺”——二阶碰撞缺口/非平凡频谱能量：
\[
F_{m+2}\,\mathsf{Col}_m-1=\frac{1}{4^m}\sum_{t\ne 0}\bigl|\widehat c_m(t)\bigr|^2,
\]
从而把 prime-like（Euler product）侧与谱/迹侧在同一张可机核验的交换图中对齐：Witt 侧把 \(\zeta=\det(I-zT)^{-1}\) 的 Euler 乘积严格从 trace 读出，而 Fourier 侧把 Fold 的 residue 叠加严格对角化，并将“偏离均匀”的强度压缩为同一个标量指纹 \(F_{m+2}\mathsf{Col}_m-1\)（等价于二阶碰撞比 \(R_m-1\)，见推论 \ref{cor:pom-renyi2-near-uniform}）。
\end{remark}

\begin{remark}[\texorpdfstring{$u$}{u} 的解释：把“不可并行 delay/串行代价”乘法化]\label{rem:zetaK-u-delay}
在本文的核/扫描算子语境中，\(\kappa(e)\) 可取为每一步不可并行的代价事件数（例如 carry 触发次数、固定延迟事件数等）。则 \eqref{eq:zetaK-euler} 表明：沿 primitive 轨道 \(\gamma\) 的总代价 \(\kappa(\gamma)\) 以乘法权重 \(u^{\kappa(\gamma)}\) 进入 primitive 因子
\[
\bigl(1-z^{|\gamma|}u^{\kappa(\gamma)}\bigr)^{-1},
\]
从而“时间侧的串行代价”在 prime-like 侧被\emph{自动转译}为 Euler 乘积的乘法指数。
这正是本文所需的可检验桥包：给定带势转移核 \(M_K(u)\)，既可在迹侧计算 \(a_n(u)=\Tr(M_K(u)^n)\)，也可经 \eqref{eq:zetaK-primitive-from-trace} 严格提取 primitive 分层谱 \(B_n(u)\)，再由 \eqref{eq:zetaK-euler} 复原 Euler 乘积。
\end{remark}

\begin{corollary}[primitive 层的代价矩：\texorpdfstring{$B_n^{(k)}(1)$}{Bn(k)(1)} 的直接读出]\label{cor:zetaK-primitive-moments}
对每个 \(n\ge 1\)，令
\(
B_n(u)=\sum_{\gamma\in\mathcal{P}_K,\ |\gamma|=n} u^{\kappa(\gamma)}
\)
如 \eqref{eq:zetaK-Bn}。
则有逐项恒等式
\[
\boxed{\;
B_n'(1)=\sum_{\gamma\in\mathcal{P}_K,\ |\gamma|=n}\kappa(\gamma),\qquad
B_n''(1)=\sum_{\gamma\in\mathcal{P}_K,\ |\gamma|=n}\kappa(\gamma)\bigl(\kappa(\gamma)-1\bigr).\;}
\]
因此在 \(u=1\) 处对 \(B_n(u)\) 求导即可得到长度 \(n\) 的 primitive 原子上“串行代价事件数”的一阶/二阶矩统计；这在 \(\kappa\) 取为 carry/delay 触发计数时给出可比较的复杂度指纹。
\end{corollary}
\begin{proof}
对每条 primitive 轨道项 \(u^{\kappa(\gamma)}\) 逐项求导并在 \(u=1\) 代入即可。证毕。
\end{proof}

\begin{remark}[以顶端溢出比特为最小势的校准基准]\label{rem:zetaK-overflow-calibration}
在 Fold 窗口层，隐藏位 \(b(\omega)\)（引理 \ref{lem:pom-one-bit}）可视为“最高权重进位是否穿透”的最小不可并行影子；
其全局计数 \(|B_m|=\lfloor 2^m/3\rfloor\) 及其无取整闭式见定理 \ref{thm:pom-hidden-bit-count} 与注 \ref{rem:pom-hidden-bit-dyadic-trunc}。
因此当 \(\kappa\) 选为更细的 carry/delay 事件计数时，可用“将 \(\kappa\) 粗化到顶端进位指示”这一投影来回收上述闭式，作为 \(u\)-势与实现管线的最小审计基准。
\end{remark}

\begin{remark}[从 toy-RH 到 RH：有限维极点格点的结构性门槛]\label{rem:finite-pole-lattice-barrier}
本文大量使用的有限维口径（邻接矩阵/有限状态核）在审计性与可计算性上有决定性优势；但在 \(s\)-平面“临界线几何”的层面，它也携带一个不可回避的结构性限制。

对任意有限维矩阵核 \(M\) 与 \(\lambda=\rho(M)\)，令
\(
\widehat{\zeta}(s)=\det(I-\lambda^{-s}M)^{-1}
\)。
命题 \ref{prop:finite-rh-triangle-dict} 给出极点的完全字典：每个非零特征值 \(\mu\) 产生一条竖直直线
\(
\Re(s)=\log|\mu|/\log\lambda
\)
上的等差格点族（公差 \(2\pi/\log\lambda\)），因此在任意有限维模型中，临界带内的“非平凡极点几何”只能由\textbf{有限条竖直直线 + 有限相位的等差梳齿}组成。
这解释了为何 40 状态核的 toy-RH（命题 \ref{prop:finite-rh-40}）虽能把极点压到 \(\Re(s)=\tfrac12\)，但临界线上的虚部结构仍是高度刚性的（由有限多个相位决定）。

因此，若希望出现更丰富的“临界线几何/统计”，需要引入极限过程（状态数增长的 sofic/SFT 逼近序列）或直接进入无穷维转移算子的 Fredholm 行列式口径；后者允许出现超出有限维“相位梳齿”的奇点结构（\cite{Ruelle1976ZetaExpanding,Ruelle1978Thermodynamic,ParryPollicott1990Zeta}）。
\end{remark}

\begin{remark}[余量扩张表示的可迭代细化：候选 \texorpdfstring{$\mathcal{L}_s$}{Ls} 与“粗分割退化为有限核”]\label{rem:remainder-map-transfer-operator-refinement}
本段把 \S\ref{subsec:online-delay} 的余量动力学
\[
R(t)=\Bigl\{\varphi R(s)+d\varphi^{-3}\Bigr\},\qquad
e=\Bigl\lfloor \varphi R(s)+d\varphi^{-3}\Bigr\rfloor\in\{0,1\},
\]
视为一个在 \(X:=[0,1)\) 上的\textbf{多分支扩张-取整系统}：对每个允许的分支对 \((d,e)\in\{0,1,2\}\times\{0,1\}\)，定义逆支（仿射收缩）
\[
\tau_{d,e}(r):=\frac{r+e-d\varphi^{-3}}{\varphi},
\]
并令 \(I_{d,e}:=\tau_{d,e}(X)\cap X\) 为其像区间（当 \((d,e)\) 不可达时该区间为空）。则对任意函数 \(f\)（例如有界变差函数或解析函数），可定义一个转移算子族
\[
(\mathcal{L}_s f)(r)
:=\sum_{(d,e)}\mathbf{1}_{I_{d,e}}(r)\,w_{d,e}(r;s)\,f\!\bigl(\tau_{d,e}(r)\bigr),
\]
其中权函数 \(w_{d,e}(r;s)\) 可取为 \(e^{-s}\)（每步长度势）、也可取为更一般的解析权（例如把 carry/delay 事件势、或 \S\ref{subsec:xi-from-self-dual} 的连续尺度混合作为 \(s\)-依赖的非局部常势），从而把 \(\det(I-\mathcal{L}_s)\) 置于 Fredholm 行列式范式。
\par
关键点在于：\textbf{粗分割时它退化回本文已用的有限核}。具体地，\S\ref{subsec:online-delay} 已给出同步核的 $10$ 个状态 \(Q\) 与边表；并且该边表可等价解释为 \(X=[0,1)\) 上的一个 Markov 分割 \(\{J_q\}_{q\in Q}\)：
对每个 \(q\in Q\)，令 \(J_q\subset[0,1)\) 为使当前余量落在该区间时“可实现状态”为 \(q\) 的取值域（端点由分支不连续点与其有限次原像生成）。
则在分段常值子空间
\(
\mathcal{C}_Q:=\mathrm{span}\{\mathbf{1}_{J_q}:q\in Q\}
\)
上，\(\mathcal{L}_s\) 的作用由一个有限矩阵完全刻画；当取 \(w_{d,e}\equiv 1\) 时，该矩阵正是同步核的邻接矩阵转置 \(B^{\mathsf T}\)（附录 \ref{app:sync-kernel-basic}），亦即
\[
\mathcal{L}_0|_{\mathcal{C}_Q}\ \cong\ B^{\mathsf T}.
\]
\par
更进一步，若以分支不连续点集合 \(\mathcal{D}_0:=\{0,1\}\cup\partial I_{d,e}\) 为初始切点，并递推细化
\[
\mathcal{D}_{m+1}:=\mathcal{D}_m\ \cup\ \bigcup_{(d,e)}\tau_{d,e}(\mathcal{D}_m),
\qquad
\Pi_m:=\text{由 }\mathcal{D}_m\text{ 诱导的区间分割},
\]
则得到一族可迭代的 Markov 细化分割 \(\Pi_m\)。在每个 \(\Pi_m\) 上取分段常值空间 \(\mathcal{C}_m\)，即可得到一串有限秩逼近
\(
\mathcal{L}_s|_{\mathcal{C}_m}
\),
其矩阵元素由“哪个逆支把某个原子映到另一个原子”逐项给出；并在 \(m=0\)（等价于同步核粗分割）时退化为上面的有限核 \(B^{\mathsf T}\)。
这提供了一个最小可执行的升级接口：保留可审计的有限矩阵层，同时允许在 \(m\to\infty\) 的极限中进入真正的无穷维 Fredholm 行列式框架。
\end{remark}

\subsubsection{柱投影与最大熵态：C*-代数口径}

与附录 \ref{app:op-algebra} 的“投影读出”语言相呼应：在 SFT 上，柱事件对应一族两两正交投影，它们生成一个交换子代数；最大熵态（Parry 测度）给出其上的规范态。更进一步，Cuntz--Krieger 代数 $\mathcal{O}_A$ 为该 SFT 提供一个非交换的普适 C*-代数封装，其对角子代数与柱投影自然同构。

\begin{remark}[柱投影的 C*-代数封装（标准事实）]\label{prop:ck-diagonal-cylinders}
令 $A$ 为不可约的 $0$--$1$ 矩阵，$X_A^+$ 为一侧 SFT。设 $\mathcal{O}_A$ 为 Cuntz--Krieger 代数（\cite{CuntzKrieger1980}），并记其标准部分等距元为 $(S_i)_{i\in\mathcal{A}}$（字母表 $\mathcal{A}$ 的索引）。对任意允许词 $\mu=\mu_1\cdots\mu_m$，记 $S_\mu:=S_{\mu_1}\cdots S_{\mu_m}$ 并定义
\[
P_\mu:=S_\mu S_\mu^*.
\]
则：
\begin{enumerate}
  \item \(P_\mu\) 族生成一个交换 C*-子代数 \(\mathcal{D}_A\subset\mathcal{O}_A\)（称为对角子代数）；
  \item 存在自然同构 \(\mathcal{D}_A\cong C(X_A^+)\)，在该同构下 \(P_\mu\) 对应于长度 \(m\) 柱集 \([\mu]\subset X_A^+\) 的指示函数；
  \item 因而任何 SFT 上的概率测度（特别是 Parry 最大熵测度）都诱导 \(\mathcal{D}_A\) 上的一个态。
\end{enumerate}
\end{remark}

\subsubsection{本文的新颖桥接点：分辨率折叠作为可审计的子代数选择}

附录 \ref{app:op-algebra-fold-subalgebra} 给出一个对本文尤为关键、但在标准 PF/热力学形式叙述中并不自动出现的结构：对任意分辨率 $m$，由全部窗口微态投影生成的有限维交换代数 $\mathcal{A}_m$，以及由折叠不变性挑出的子代数 $\mathcal{B}_m$。该子代数选择正是“分辨率压缩/稳定类型化”的算子代数版本：它把可观测量限制为对折叠纤维常值的部分，从而把微态层的巨大自由度压缩为 $|X_m|=F_{m+2}$ 个稳定类型坐标。\par
在此意义下，本文的算子论新信息不是重新证明 PF/Parry/$\zeta$ 的经典结论，而是给出一条\textbf{从扫描投影读出到语法层算子接口的可计算桥接链}：\[
\text{扫描-投影二元读出}\ \Rightarrow\ \text{窗口微态投影代数 }\mathcal{A}_m\ \Rightarrow\ \text{折叠不变子代数 }\mathcal{B}_m\ \Rightarrow\ \text{语法层（SFT/sofic）与 }\zeta\text{ 接口}.
\]

\begin{remark}[折叠专属的新接口：纤维条件期望与多重度]\label{rem:fold-new-invariants}
在标准的 SFT/\(\zeta\) 叙述中，“对角子代数/柱投影”通常是给定语法后自然出现的；而在本文生成链中，\(\mathcal{B}_m\subset\mathcal{A}_m\) 的出现是由折叠 \(\Fold_m\) 强制的 coarse-graining。更进一步，附录 \ref{prop:fold-conditional-expectation} 给出一个完全显式的条件期望
\[
E_m:\mathcal{A}_m\to\mathcal{B}_m,\qquad
E_m(P_a)=\frac{1}{d_m(\Fold_m(a))}\widetilde P_{\Fold_m(a)},
\]
其系数由折叠纤维多重度 \(d_m(x)=|\Fold_m^{-1}(x)|\) 给出。该 \(E_m\) 可视为“信息丢失接口”：它把微态层可观测量压缩为对折叠纤维常值的部分，并将多重度作为一个可审计输入。

在核/同步图层面，本文进一步把“折叠/归一化机制”编码为有限图移位，从而引入一组可比较的不变量族：例如 Abel/Hadamard 有限部常数 \(\log\mathfrak{M}\)、orbit--Mertens 常数 \(\mathsf{M}\)、以及 Dirichlet--Mertens 常数张量与最坏扭曲谱半径比 \(\rho/\lambda\)。这些对象在形式上仍由 \(\zeta=1/\det(I-zA)\) 的解析结构导出，但其\emph{输入矩阵}来自归一化核的同步图（而非仅来自语法本身），因此可作为“折叠专属的谱指纹”用于跨核对比。
\end{remark}
